% ============================================================
%  CHƯƠNG I: TỔNG QUAN
% ============================================================
\chapter{TỔNG QUAN}
\label{chap:tong_quan}

% ────────────────────────────────────────────────────────────
\section{Giới thiệu đề tài}
\label{sec:gioi_thieu}

\subsection{Bối cảnh}
\label{subsec:boi_canh}

Trong những năm gần đây, thị trường cho thuê nhà ngắn hạn tại Việt Nam và trên toàn cầu đã chứng kiến
sự tăng trưởng vượt bậc nhờ sự phổ biến của các nền tảng trực tuyến như Airbnb, Booking.com hay
Agoda. Mô hình này không chỉ tạo ra nguồn thu nhập thụ động cho chủ nhà mà còn cung cấp cho
khách du lịch những lựa chọn đa dạng, linh hoạt hơn so với khách sạn truyền thống.

Tuy nhiên, đằng sau giao diện người dùng đơn giản là những bài toán kỹ thuật phức tạp mà các
hệ thống này phải giải quyết:

\begin{itemize}
    \item \textbf{Tính toàn vẹn dữ liệu:} Một bất động sản không thể bị đặt trùng bởi hai khách
    hàng cùng lúc — đây là bài toán xử lý đồng thời (concurrency control) trong hệ thống phân tán.

    \item \textbf{Khả năng tìm kiếm theo địa lý:} Khách cần tìm phòng trong bán kính nhất định
    so với điểm du lịch — đòi hỏi hệ thống tìm kiếm địa không gian (geo-spatial search) hiệu năng cao.

    \item \textbf{Đồng bộ dữ liệu thời gian thực:} Khi chủ nhà cập nhật thông tin, kết quả tìm
    kiếm của khách phải phản ánh ngay lập tức mà không cần truy vấn trực tiếp vào cơ sở dữ liệu
    giao dịch.

    \item \textbf{Giao tiếp theo ngữ cảnh:} Hệ thống chat không đơn thuần là nhắn tin P2P mà phải
    được gắn chặt với ngữ cảnh nghiệp vụ (một bất động sản hoặc một đơn đặt phòng cụ thể).

    \item \textbf{Xử lý thanh toán an toàn:} Luồng thanh toán phải đảm bảo tính nguyên tử
    (atomicity) và không block hệ thống — kết quả thanh toán phải được xử lý bất đồng bộ.
\end{itemize}

Đề tài \textbf{AirVnV} ra đời từ nhu cầu nghiên cứu và xây dựng một nền tảng cho thuê nhà ngắn
hạn hiện đại, áp dụng kiến trúc phân tán để giải quyết các thách thức nêu trên trong điều kiện
thực tế. Đây không phải là bài toán học thuật đơn thuần mà là một bài toán kỹ thuật phần mềm
có tính ứng dụng cao.

\subsection{Mục tiêu}
\label{subsec:muc_tieu}

Mục tiêu của đồ án là xây dựng một nền tảng đặt phòng trực tuyến hoàn chỉnh — từ backend đến
frontend — với các chức năng cốt lõi gồm: quản lý bất động sản, tìm kiếm theo địa lý, đặt phòng,
thanh toán và chat theo ngữ cảnh.

Bên cạnh việc hoàn thiện sản phẩm về mặt chức năng, nhóm tập trung vào việc giải quyết
các bài toán kỹ thuật thực tế trong hệ thống phân tán:

\begin{itemize}
    \item Áp dụng kiến trúc Microservices — mỗi nghiệp vụ là một service độc lập, có thể
    phát triển và triển khai riêng lẻ.

    \item Đảm bảo tính nhất quán dữ liệu giữa các service mà không dùng distributed transaction,
    thông qua Transactional Outbox Pattern và Change Data Capture (CDC).

    \item Xây dựng hệ thống tìm kiếm địa lý hiệu năng cao bằng Elasticsearch — kết quả tìm kiếm
    không truy vấn trực tiếp vào database giao dịch.

    \item Xây dựng hệ thống chat real-time gắn với ngữ cảnh đặt phòng, tự động phát sinh
    thông báo hệ thống khi có sự kiện nghiệp vụ.

    \item Thiết kế luồng thanh toán bất đồng bộ — kết quả thanh toán không block luồng chính
    mà được xử lý qua message broker.
\end{itemize}

\subsubsection*{Phạm vi}

Đồ án bao gồm ba nhóm người dùng: Khách (Guest), Chủ nhà (Host) và Quản trị viên (Admin).
Ứng dụng mobile native nằm ngoài phạm vi của đồ án này.


% ────────────────────────────────────────────────────────────
\section{Cơ sở lý thuyết}
\label{sec:co_so_ly_thuyet}

Để giải quyết các bài toán đặt ra, đồ án áp dụng một tập hợp các công nghệ và kiến trúc hiện đại.
Phần này trình bày nền tảng lý thuyết của từng thành phần và lý do lựa chọn.

% ── 2.1 Microservices Architecture ──────────────────────────
\subsection{Kiến trúc Microservices}
\label{subsec:microservices}

Kiến trúc Microservices phân chia hệ thống thành các service nhỏ, độc lập, mỗi service chịu
trách nhiệm về một miền nghiệp vụ duy nhất và có cơ sở dữ liệu riêng (Database-per-Service).
Ngược lại với kiến trúc Monolith, Microservices cho phép:

\begin{itemize}
    \item Triển khai độc lập từng service mà không cần build lại toàn bộ hệ thống.
    \item Scale theo chiều ngang (horizontal scaling) có chọn lọc — chỉ scale service đang chịu tải cao.
    \item Cách ly lỗi (fault isolation): lỗi ở một service không làm sập toàn bộ hệ thống.
\end{itemize}

Trong AirVnV, hệ thống được phân rã thành 7 service: \textit{UserService, PropertyService,
SearchService, BookingService, PaymentService, ChatService} và \textit{Gateway}.

% ── 2.2 API Gateway Pattern ─────────────────────────────────
\subsection{API Gateway Pattern}
\label{subsec:api_gateway}

API Gateway là một điểm vào duy nhất (single entry point) cho tất cả các yêu cầu từ client.
Gateway đảm nhận các trách nhiệm:

\begin{itemize}
    \item \textbf{Routing:} Định tuyến request đến đúng service.
    \item \textbf{Authentication:} Xác thực JWT token trước khi forward request.
    \item \textbf{Load Balancing:} Phân phối tải đến các instance của cùng một service.
\end{itemize}

Đồ án sử dụng \textbf{YARP} (Yet Another Reverse Proxy) — một thư viện reverse proxy hiệu năng cao
được phát triển bởi Microsoft, tích hợp tốt với .NET ecosystem và .NET Aspire.

% ── 2.3 CQRS và Mediator Pattern ────────────────────────────
\subsection{CQRS và Mediator Pattern}
\label{subsec:cqrs_mediator}

\textbf{CQRS (Command Query Responsibility Segregation)} là pattern tách biệt hoàn toàn luồng ghi dữ liệu (Command) và luồng đọc dữ liệu (Query).
Mỗi thao tác nghiệp vụ được biểu diễn bằng một object (Command hoặc Query) và xử lý bởi một Handler tương ứng.

\begin{itemize}
    \item \textbf{Command:} Thay đổi trạng thái hệ thống. Ví dụ: \texttt{CreateBookingCommand},
    \texttt{PublishPropertyCommand}.
    \item \textbf{Query:} Chỉ đọc dữ liệu, không có side-effect. Ví dụ: \texttt{GetPropertyQuery},
    \texttt{SearchPropertiesQuery}.
\end{itemize}

Để triển khai CQRS một cách hiệu quả và giảm thiểu sự phụ thuộc chéo (coupling) giữa các class, đồ án áp dụng \textbf{Mediator Pattern} (thông qua thư viện \texttt{Mediator.SourceGenerator}).
Mediator Pattern đóng vai trò là một "người điều phối" trung tâm. Thay vì các Endpoint/Controller phải tiêm (inject) hàng loạt các Service/Repository gây ra tình trạng "God object",
chúng chỉ cần gửi (Send) một \texttt{Request} tới Mediator. Mediator sẽ dựa vào kiểu dữ liệu của Request để tự động định tuyến tới đúng \texttt{Handler} duy nhất chịu trách nhiệm xử lý.

Sự kết hợp giữa CQRS, Mediator Pattern và \textbf{FastEndpoints} giúp hệ thống hiện thực hóa mô hình kiến trúc cắt dọc (\textbf{Vertical Slice Architecture}).
Theo đó, mỗi Use Case nghiệp vụ được gom cụm độc lập trong một thư mục duy nhất (gồm \texttt{Endpoint.cs, Request.cs, Response.cs, Handler.cs}),
tối đa hóa tính gắn kết nội bộ (high cohesion) và hạn chế tối đa rủi ro khi thay đổi mã nguồn.

% ── 2.4 Event-Driven Architecture ───────────────────────────
\subsection{Kiến trúc Hướng sự kiện (Event-Driven Architecture)}
\label{subsec:eda}

Event-Driven Architecture (EDA) là mô hình các service giao tiếp với nhau thông qua việc
phát và tiêu thụ sự kiện (events) bất đồng bộ thay vì gọi trực tiếp (synchronous HTTP).

Trong AirVnV, hai message broker được sử dụng với vai trò khác nhau:

\begin{itemize}
    \item \textbf{RabbitMQ:} Dùng cho Domain Events nội bộ — các sự kiện nghiệp vụ ngắn hạn
    giữa các service (ví dụ: \texttt{BookingConfirmedEvent} kích hoạt ChatService inject
    system message).

    \item \textbf{Apache Kafka:} Dùng cho Data Streaming — luồng dữ liệu lớn, bền vững, hỗ
    trợ replay. Cụ thể là luồng CDC từ PostgreSQL đến Elasticsearch.
\end{itemize}

% ── 2.5 Transactional Outbox Pattern ────────────────────────
\subsection{Transactional Outbox Pattern}
\label{subsec:outbox}

Một thách thức cốt lõi của Microservices là: làm sao đảm bảo rằng khi một nghiệp vụ thành công
(ví dụ: lưu Booking vào DB), sự kiện tương ứng (ví dụ: \texttt{BookingCreatedEvent}) luôn
được publish đến message broker — kể cả khi ứng dụng crash ngay sau khi lưu DB?

Outbox Pattern giải quyết vấn đề này bằng cách lưu event vào một bảng \texttt{OutboxMessage}
trong \textit{cùng một transaction} với nghiệp vụ chính. Một background worker riêng biệt
sẽ đọc bảng Outbox và publish event đến broker, đảm bảo tính nguyên tử (atomicity) mà không
cần distributed transaction.

% ── 2.6 Change Data Capture ─────────────────────────────────
\subsection{Change Data Capture (CDC) với Debezium}
\label{subsec:cdc}

CDC là kỹ thuật theo dõi các thay đổi trong cơ sở dữ liệu (INSERT, UPDATE, DELETE) ở
mức transaction log và phát chúng ra ngoài như các sự kiện.

Đồ án sử dụng \textbf{Debezium} — một nền tảng CDC mã nguồn mở — để đọc Write-Ahead Log (WAL)
của PostgreSQL và publish các thay đổi vào Kafka. SearchService tiêu thụ luồng Kafka này để
cập nhật Elasticsearch, đảm bảo read-model luôn đồng bộ gần thời gian thực với dữ liệu giao dịch.

% ── 2.7 Geo-Spatial Search ──────────────────────────────────
\subsection{Tìm kiếm và Chỉ mục với Elasticsearch}
\label{subsec:geo_search}

Elasticsearch là một công cụ tìm kiếm và phân tích phân tán (Distributed Search Engine) dựa trên Apache Lucene.
Thay vì lưu dữ liệu dạng bảng như RDBMS và sử dụng phép toán \texttt{LIKE} vốn cực kỳ chậm chạp trên tập dữ liệu lớn,
Elasticsearch sử dụng cấu trúc \textbf{Inverted Index} (Chỉ mục ngược). Cấu trúc này ánh xạ từng từ khóa (token)
ngược lại với các tài liệu (documents) chứa từ khóa đó, giúp tốc độ tìm kiếm văn bản toàn văn (Full-text Search)
nhanh gấp hàng trăm lần SQL.

Bên cạnh đó, Elasticsearch hỗ trợ mạnh mẽ kiểu dữ liệu \texttt{geo\_point}, cho phép thực hiện các truy vấn
không gian (Geo-Spatial) phức tạp như:
\begin{itemize}
    \item \textbf{Geo Distance Query:} Tìm tất cả bất động sản trong bán kính $r$ km từ một tọa độ GPS $(lat, lon)$.
    \item \textbf{Geo Bounding Box:} Lọc các bất động sản nằm lọt thỏm trong một vùng khung chữ nhật trên bản đồ.
\end{itemize}
Thông qua thư viện client \texttt{Aspire.Elastic.Clients.Elasticsearch}, \texttt{SearchService} thao tác trực tiếp với
Elasticsearch dưới dạng Read-Model. Việc tách biệt hoàn toàn cơ sở dữ liệu đọc (Elasticsearch) và ghi (PostgreSQL)
giúp hiệu năng tìm kiếm đạt được tốc độ cực cao, $O(1)$, bất chấp sự gia tăng của lượng dữ liệu.

% ── 2.8 .NET Aspire ─────────────────────────────────────────
\subsection{.NET Aspire — Service Discovery và Orchestration}
\label{subsec:aspire}

.NET Aspire là nền tảng orchestration của Microsoft dành cho ứng dụng phân tán. Thay vì
hardcode connection string và URL giữa các service, Aspire cung cấp:

\begin{itemize}
    \item \textbf{Service Discovery:} Các service tìm thấy nhau bằng tên logic
    (ví dụ: \texttt{http://property-service/...}) thay vì IP/port cứng.
    \item \textbf{Infrastructure Provisioning:} Aspire tự động khởi động các container
    (PostgreSQL, Redis, Kafka, RabbitMQ, Elasticsearch) khi chạy môi trường phát triển.
    \item \textbf{Observability Dashboard:} Giao diện theo dõi logs, metrics và traces
    tập trung cho toàn bộ hệ thống.
\end{itemize}

% ── 2.9 Hệ quản trị CSDL PostgreSQL & ORM ───────────────────
\subsection{Hệ quản trị Cơ sở dữ liệu PostgreSQL và ORM}
\label{subsec:db_orm}

\textbf{PostgreSQL} được lựa chọn làm hệ quản trị cơ sở dữ liệu quan hệ (RDBMS) chính cho các Microservices.
Lý do chính bao gồm: tính tuân thủ ACID nghiêm ngặt giúp đảm bảo toàn vẹn dữ liệu tài chính (Booking, Payment),
và khả năng hỗ trợ kiểu dữ liệu \texttt{JSONB} mạnh mẽ. \texttt{JSONB} cho phép lưu trữ các trường dữ liệu bán cấu trúc
(semi-structured data) như danh sách tiện ích (amenities) hoặc bộ lọc tìm kiếm một cách linh hoạt mà không cần
chuẩn hóa quá mức (over-normalization).

Để giao tiếp với PostgreSQL, đồ án sử dụng \textbf{Entity Framework Core} (EF Core) thông qua thư viện \texttt{Npgsql.EntityFrameworkCore.PostgreSQL}.
EF Core là một Object-Relational Mapper (ORM) mạnh mẽ, cung cấp cơ chế Code-First Migrations, giúp
tự động hóa việc tạo và cập nhật schema cơ sở dữ liệu. Để đảm bảo nguyên tắc \textit{Database-per-service},
mỗi Microservice duy trì một lịch sử Migration độc lập và tự động apply vào Logical Database của riêng nó
ngay khi khởi động.

% ── 2.10 Hệ thống Real-time với SignalR & Redis Backplane ───
\subsection{Hệ thống Thời gian thực với SignalR và Redis Backplane}
\label{subsec:signalr_redis}

Chức năng nhắn tin (Chat) và thông báo (Notification) đòi hỏi kết nối hai chiều liên tục giữa Server và Client
với độ trễ cực thấp. Để đáp ứng nhu cầu này, đồ án sử dụng \textbf{SignalR} (\texttt{Microsoft.AspNetCore.SignalR}).
SignalR tự động quản lý các kết nối WebSocket (và fallback xuống Server-Sent Events hoặc Long Polling nếu cần),
cho phép máy chủ đẩy dữ liệu (push) trực tiếp đến các client đang kết nối.

Tuy nhiên, trong môi trường phân tán (Docker/Kubernetes), \texttt{ChatService} có thể được nhân bản (scale-out) thành nhiều node.
Nếu User A kết nối tới Node 1 và User B kết nối tới Node 2, Node 1 sẽ không biết cách gửi tin nhắn trực tiếp cho User B.
Để giải quyết rào cản này, \textbf{Redis Backplane} (\texttt{StackExchange.Redis}) được tích hợp.
Mọi tin nhắn SignalR gửi ra sẽ được đẩy vào kênh Redis Pub/Sub; sau đó Redis sẽ broadcast tin nhắn này tới tất cả
các node của \texttt{ChatService}, đảm bảo tin nhắn luôn đến đúng người nhận bất kể họ đang duy trì socket với node nào.

% ── 2.11 Bảo mật, Định tuyến và Ánh xạ Dữ liệu ──────────────
\subsection{Bảo mật, Thiết kế API và Ánh xạ dữ liệu}
\label{subsec:security_mapping}

\begin{itemize}
    \item \textbf{Thiết kế API siêu tốc (FastEndpoints):} Thay vì sử dụng mô hình MVC/Controllers truyền thống vốn
    cồng kềnh, hệ thống sử dụng thư viện \texttt{FastEndpoints}. Thư viện này triển khai kiến trúc REPR
    (Request-Endpoint-Response), giúp nhóm (group) tất cả logic nghiệp vụ, validation và routing của một API vào một file duy nhất,
    tối ưu tốc độ thực thi và cực kỳ dễ dàng bảo trì.
    
    \item \textbf{Bảo mật \& Xác thực (Security):} Hệ thống kết hợp \textbf{JSON Web Token (JWT)} để xác thực không
    trạng thái (stateless authentication) qua YARP Gateway. Mật khẩu người dùng được băm (hash) an toàn bằng thuật toán bcrypt qua thư viện
    \texttt{BCrypt.Net-Next}. Đồng thời, đồ án tích hợp hệ thống đăng nhập một chạm (SSO) qua Google nhờ các thư viện
    \texttt{FirebaseAdmin} và \texttt{Google.Apis.Auth}.
    
    \item \textbf{Ánh xạ dữ liệu (Mapster):} Trong kiến trúc sạch, dữ liệu phải được ánh xạ liên tục giữa Entity và DTO. 
    Thay vì dùng AutoMapper (dựa trên Reflection làm suy giảm hiệu năng), hệ thống sử dụng \texttt{Mapster}.
    Mapster sinh mã trực tiếp lúc biên dịch (code generation), giúp tốc độ chuyển đổi dữ liệu đạt mức tối đa.
\end{itemize}

% ── 2.12 Tích hợp Thanh toán Ngoại vi ───────────────────────
\subsection{Tích hợp Cổng thanh toán Ngoại vi (Payment Gateway)}
\label{subsec:payment_gateway}

Luồng thanh toán trong AirVnV không lưu trữ trực tiếp thông tin thẻ tín dụng nhằm giảm tải rủi ro pháp lý và
đảm bảo tuân thủ tiêu chuẩn PCI-DSS. Thay vào đó, \texttt{PaymentService} tích hợp với nền tảng \textbf{Stripe} thông qua thư viện \texttt{Stripe.net} (hoặc VNPay).

Khi khách hàng khởi tạo thanh toán, hệ thống sẽ sinh ra một Payment Session và điều hướng người dùng sang giao diện an toàn của đối tác.
Mấu chốt của kiến trúc thanh toán phân tán là cơ chế \textbf{Webhook}: Hệ thống thanh toán sẽ gọi ngược (callback)
về hệ thống AirVnV thông qua một API đặc biệt được bảo mật bằng chữ ký số (Webhook Signature) để báo cáo kết quả giao dịch.
Nhờ luồng xử lý này, hệ thống backend không bao giờ bị "treo" (block) để chờ kết quả giao dịch, đảm bảo tính bất đồng bộ 
và chống rớt đơn hàng hiệu quả kể cả khi kết nối mạng từ client bị ngắt ngang.
